\section{Circular orbits}
% #physics

How fast does a satellite go, and how long does one turn take? Only gravity holds it, so the two are tied together.

A satellite of mass $m$ circles a planet of mass $M$ at a distance $r$, and gravity keeps it turning:
\begin{equation}
  \frac{G M m}{r^2} = \frac{m v^2}{r}
  \qquad\text{so}\qquad
  v = \sqrt{\frac{G M}{r}} .
\end{equation}

\begin{theorem}[Kepler's third law]
The square of the period grows as the cube of the distance:
\[ T^2 = \frac{4\pi^2}{G M}\, r^3 . \]
\end{theorem}

\begin{proof}
One turn is $2\pi r$ long, so $T = 2\pi r / v$. Square it, and put in $v^2 = G M / r$.
\end{proof}

\begin{tabular}{lrr}
  & Distance (km) & Period \\
  \hline
  Space station & 6,771 & 92 minutes \\
  GPS satellite & 26,560 & 12 hours \\
  The Moon & 384,400 & 27.3 days \\
\end{tabular}
